\songtitle{\textjapanese{四つ目の扉} (Yottsume no Tobira)}

\tag*{2026}\tag{anime-pop}\tag{japanese-lyrics}\tag{multilingual}\tag{african-empire}\tag{monukae}\tag{mombuto}\tag{gladiel}
% -> stories.tex (African Empire -- Monukae)
\begin{notes}
Lyrics by Claude based on the \emph{Monukae} saga within the \emph{African Empire} alternate-history universe by Carlos Thompson, written as an anime theme song: a duet of Mombuto and Gladiel facing the emperor's two-month deadline, three bad doors and the search for a fourth. Japanese with Irish, Latin and Greek lines. Retrieved from Wyomee Dank's Suno page (V6-MINI, 14 November 2026).
\end{notes}
\begin{style}
Cinematic chamber-orchestral anime pop; kora motif and piano ostinato thread through cello and low strings, harpsichord, frame drum and ticking clock percussion; low organ swells before a hushed choir enters on Latin and Greek chants in the bridge, returning for the key-raised final chorus; suspended final chord; male baritone and female mezzo duet in Japanese, restrained, whispered, almost spoken in the verses, crooning legato phrasing and whispered ad-libs between lines, then harmonizing together; tense, dramatic, wide-stereo but restrained, minor-key slow burn at 75 BPM.
\end{style}
\begin{lyrics}
\songpart{Intro}
\annotation{Instrumental: slowed kora motif, ticking clock, low cello}
\annotation{Whispered female vocal}\\
Ní dhéanfaidh mé dearmad

\songpart{Verse 1}
\annotation{Male vocal, restrained baritone, almost spoken}\\
\begin{japanese}
笑顔の裏で 数を数える\\
敵と味方の 線を引く\\
父の影が 長すぎて\\
歩幅が 測れない
\end{japanese}

\annotation{Female vocal, mezzo, answering}\\
\begin{japanese}
憎んだ名を 呼んでみる\\
舌の上で 苦い味\\
憎しみは 鎖にも\\
鍵にも なると 知った
\end{japanese}

\songpart{Pre-Chorus}
\annotation{Duet, strings climbing, ticking clock louder}\\
\begin{japanese}
にかげつの 期限\\
三つの扉 どれを選んでも\\
何かを 失う\\
奪われる前に 差し出す
\end{japanese}

\songpart{Chorus}
\annotation{Duet, harmonies join for the first time, epic but controlled}\\
\begin{japanese}
譲るふりして 全てを守る\\
膝を折る 手の裏で\\
三つの扉の 外にある\\
四つ目の扉を 探している
\end{japanese}

\songpart{Verse 2}
\annotation{Female vocal}\\
\begin{japanese}
偽りの指輪さえ\\
誓いに 変わっていく\\
\end{japanese}
\annotation{Male vocal}\\
\begin{japanese}
異国の神の 前で\\
頭を垂れて 笑う\\
\end{japanese}
\annotation{Duet, low}\\
\begin{japanese}
信じてなくても 祈る\\
扉の鍵に なるなら
\end{japanese}

\songpart{Pre-Chorus}
\annotation{Low organ, hushed choir, Latin chant}\\
Ego te baptizo
\annotation{Male vocal, low}\\
\begin{japanese}
水に沈む 古い名\\
新しい名で 扉を叩く
\end{japanese}

\songpart{Chorus 2}
\annotation{Duet, fuller strings}\\
\begin{japanese}
譲るふりして 全てを守る\\
その嘘を 誰が責める\\
三つの扉の 外にある\\
四つ目の扉を 開けてゆく
\end{japanese}

\songpart{Bridge}
\annotation{Only ticking clock and low piano, hushed whispered voices}\\
\begin{japanese}
疑いは 静かに 育つ\\
玉座の目が こちらを見る\\
証拠は ない\\
だから 怖いんだ\\
\end{japanese}
\annotation{Choir, hushed, Greek chant}\\
Kyrie eleison
\annotation{Silence for one beat}

\songpart{Final Chorus}
\annotation{Duet in unison and harmony, full orchestra and choir, key change up}\\
\begin{japanese}
譲るふりして 全てを守る\\
膝を折る 手の裏で\\
三つの扉の 外にある\\
四つ目の扉は もう 開いている
\end{japanese}

\songpart{Outro}
\annotation{Kora motif alone, ticking resumes then stops, final suspended chord, unresolved}
\annotation{Male vocal, whispered}\\
\begin{japanese}
筆が 静かに 動く
\end{japanese}

\annotation{End}
\end{lyrics}

